% Generated draft: requires mathematical and source review.
\documentclass[11pt]{article}
\usepackage[margin=25mm]{geometry}
\usepackage{fontspec}
\setmainfont{Latin Modern Roman}
\usepackage{amsmath,amssymb,mathtools,unicode-math}
\setmathfont{Latin Modern Math}
\usepackage{enumitem,needspace,xurl,hyperref}
\setlength{\parindent}{0pt}
\setlength{\parskip}{5pt}
\setlength{\emergencystretch}{3em}
\newcommand{\N}{\mathbb{N}}
\newcommand{\Z}{\mathbb{Z}}
\newcommand{\Q}{\mathbb{Q}}
\newcommand{\R}{\mathbb{R}}
\newcommand{\C}{\mathbb{C}}
\newcommand{\F}{\mathbb{F}}
\newcommand{\A}{\mathbb{A}}
\newcommand{\PP}{\mathbb{P}}
\newcommand{\E}{\mathbb{E}}
\newcommand{\eps}{\varepsilon}
\newcommand{\dd}{\,\mathrm{d}}
\begin{document}
\section*{UnsolvedMath ID: 800425}
\textbf{Source problem code:} catalogue entry
\subsection{Definitions and mathematical statement}
Let \(d\ge 2\) be a fixed integer, and let \(q\) range over prime powers. Write \(\F_q\) for the field with \(q\) elements. For a field \(K\), define
\[
\operatorname{Poly}_d^1(K)=\{f\in K[x]:\deg f=d\},
\]
and, as in the supplied record,
\[
\operatorname{End}_d^1(K)=\left\{\frac{A}{B}\in K(x):\gcd(A,B)=1,\ \max\{\deg A,\deg B\}=d\right\}.
\]
The fractions in this definition represent rational functions, so different numerator--denominator representations of the same function do not give different elements.

The supplied excerpt uses the predicate ``\((f,0)\) is stable over \(\F_q\)'' without defining it. Retain this predicate as \(\mathsf{Stable}_{\F_q}(f,0)\); its mathematical meaning must be recovered from the source before this becomes a complete standalone statement. With that qualification, set
\[
\operatorname{Poly}_d^1(\F_q)^{\mathrm{stab}}
=\{f\in\operatorname{Poly}_d^1(\F_q):\mathsf{Stable}_{\F_q}(f,0)\}.
\]
The source also writes
\[
\operatorname{Rat}_d^1(\F_q)^{\mathrm{stab}}
=\{f\in\operatorname{Rat}_d^1(\F_q):\mathsf{Stable}_{\F_q}(f,0)\},
\]
but does not define \(\operatorname{Rat}_d^1\), despite defining \(\operatorname{End}_d^1\). If \(\operatorname{Rat}_d^1(K)\) is intended to mean \(\operatorname{End}_d^1(K)\), the rational-function counting problem concerns all degree-\(d\) rational functions satisfying the same stability predicate. This identification is conditional and requires source verification.

\textbf{Question.} For each fixed \(d\ge2\), determine the asymptotic growth, as \(q\to\infty\) through prime powers, of
\[
\#\operatorname{Poly}_d^1(\F_q)^{\mathrm{stab}}
\quad\text{and}\quad
\#\operatorname{Rat}_d^1(\F_q)^{\mathrm{stab}},
\]
subject to resolving the definitions above. These are counts of maps themselves; no quotient by conjugacy is specified.

The source additionally asks analogous counting questions with stability replaced by eventual stability, and for maps of \(\A^N\) or \(\PP^N\) with \(N\ge2\). The excerpt does not define eventual stability or specify the higher-dimensional families, degree restrictions, marked data, or corresponding stability predicates. These extensions therefore remain requests for analogous formulations, rather than fully specified counting problems. No particular asymptotic formula or growth rate is conjectured in the supplied excerpt.
\subsection{Short English statement}
For a fixed degree \(d\ge2\), count the polynomials and rational functions over \(\F_q\) for which the pair \((f,0)\) is stable, and determine how those counts grow as \(q\) increases. The supplied excerpt leaves stability undefined and uses an undefined rational-map notation. It also requests eventual-stability and higher-dimensional analogues without enough details to formulate them uniquely.
\subsection{Sources}
\begin{itemize}
\item Local UnsolvedMath registry, permanent ID \(800425\), exact source problem code \texttt{catalogue entry}, title \emph{Open Problems in Arithmetic Dynamics --- Question 48}. The supplied registry labels the problem \texttt{open}; this is a supplied status label, not an independently verified mathematical or literature assessment.
\item Benedetto, DeMarco, Ingram, Jones et al., \emph{Current Trends and Open Problems in Arithmetic Dynamics} (2018), Question 48. Supplied source URL: \url{https://arxiv.org/abs/1806.04980}. The record attributes the problem to these authors and gives proposed year 2018.
\item Supplied provenance: extraction method \texttt{arxiv\_latex}; access date \texttt{2026-10-16}; upstream revision \texttt{d371ee917ac9647fbb06cd6fc529deec6f3dda82}. The access date and provenance are supplied metadata; no independent source visit or verification is asserted.
\end{itemize}
\end{document}
